La epilepsia es uno de los trastornos neurológicos más prevalentes a nivel mundial, afectando a aproximadamente 50 millones de personas. En Guatemala, se estima que alrededor de 325,000 individuos padecen esta condición. A pesar de la existencia de diversos servicios de neurología, muchos pacientes requieren estudios especializados que se concentran en centros de referencia como el Centro de Epilepsia y Neurocirugía Funcional HUMANA. En estos centros, uno de los principales retos para el diagnóstico oportuno es la alta carga de trabajo que representa para los especialistas la evaluación visual exhaustiva de estudios multicanal como los electroencefalogramas (EEG). 

En este contexto, se propuso la aplicación de técnicas de inteligencia artificial con enfoque híbrido, a nivel morfológico y modal, como herramienta de apoyo al diagnóstico. El enfoque híbrido se orienta a la integración estructurada de distintas estrategias de modelado, procesamiento de señales, arquitecturas de aprendizaje profundo y modelos clásicos de aprendizaje automático, cuya selección se definió según el desempeño observado durante la fase de validación. La metodología contempla la implementación y evaluación de estas configuraciones para la identificación automática de segmentos de interés en las señales bioeléctricas.

La validación del sistema se realizó mediante métricas estándar de clasificación y un análisis comparativo frente a los modelos previamente desarrollados en esta línea de investigación de la Universidad del Valle de Guatemala. Finalmente, se actualizó la herramienta integrando el modelo en el software existente con el propósito de mejorar su interpretabilidad, eficiencia computacional y aplicabilidad en entornos clínicos y académicos.